\section{Concluding Remarks}
\vspace{-0.2cm}
This work addresses the critical challenge of achieving strong DP in FL for on-device LMs. We have successfully extended the BLT mechanism to multi-participation scenarios and integrated it into the DP-FTRL framework. Our BLT-DP-FTRL algorithm demonstrates superior privacy-utility trade-offs compared to the widely-used \treeagg mechanism while maintaining its ease of use. Furthermore, it rivals the state-of-the-art \bandmf mechanism in performance, yet without the associated complexities and high memory costs. Through extensive empirical evaluations on both a benchmark dataset and real-world production tasks, we have showcased the practicality and effectiveness of BLT-DP-FTRL, paving the way for its broader adoption. 

The empirical results in this paper primarily focus on the cross-device FL setting where privacy amplification by sampling is challenging in practice. The discussions (e.g. \cref{tab:mechanism_summary}) can also be applied to centralized setting for user-level DP or example-level DP. In centralized setting, \bandmf~\citep{choquette2023amplified} with amplification can achieve better privacy-utility trade-off measured by \rmsloss among the mentioned mechanisms, when number of rounds $n$ and model dimension $m$ is not too large for optimizing and applying the mechanism. When $n$ and $m$ are large, \blt and \bandtoep~\citep{mckenna2024scaling} (similarly, \bandfhu~\citep{kalinin2024banded}) can both be applied, where \blt has less optimization cost for very large $n$ (shown in \cref{fig:faster}), while \bandtoep can apply the known amplification by sampling mechanism. In conclusion, BLT-DP-FTRL algorithm is recommended for cross-device federated learning~\citep{xu2023federated,kairouz2019advances}, and is also useful for DP in datacenter~\citep{ponomareva2023dpfy,xu2022learning,charles2024fine} and federated learning in trusted execution environments (TEEs)~\citep{daly2024federated,eichner2024confidential}.